\section{Security}

\subsection{Ideal Functionality $\F{\ATS}$.} $\F{\ATS}$ receives inputs from
the system participants and simulator $\simu$ and correctly generate the
corresponding outputs, which are then transmitted to the parties authorised to
access them.  $\F{\ATS}$ exposes five interfaces: register, issue, verify.
$\simu$ and honest system participants submit their inputs through these
interfaces.

\paragraph{Register.} This is called by a user to register into the system.  If
the registration authority accepts, and the user has never been registered
before, the user is successfully registered.

\paragraph{Issuance.} This is invoked when a user wants to get a credential
that certifies value $\atv{}$ of some attribute $\att{}$.  If the authorised
issuer accepts the issuance request, the user obtains a credential that they
can use in subsequent presentations.

During a call, $\simu$ learns the identity of the user. During a call to issue,
$\simu$ additionally learns the value of the user's attribute.

\paragraph{Aggregation.} This is invoked by a user to aggregate the user's
credentials for attribute values $\vec{\atv{}}$ in a way that enables them to
present the aggregate credential later.  Aggregation is a local operation, and
so $\simu$ does not learn any information during the call.  A successful call
results in the user receiving a random string $\signature$ representing the
aggregate credentials.

\paragraph{Presentation.} Present is queried by a user to produce a
presentation of attribute values ($\vec{\atv{}}$).  the user 's request
consists of triple $(\vec{\atv{}}, \signature, \rand)$, with $\signature$ the
aggregation of the credentials certifying $\vec{\atv{}}$ and $\rand$ a random
number that identifies this presentation request.  If the user has previously
successfully called aggregate with $\vec{\atv{}}$, then $\F{\ATS}$ accepts the
query. Otherwise, it checks if the user actually holds the $\vec{\atv{}}$.
This captures the fact that even if a corrupt the user does not call aggregate
before calling present, $\F{\ATS}$ may still accept the presentation request as
long as the user has previously received the relevant credentials.  If
$\F{\ATS}$ accepts, then it generates a random string $\proof$; the randomness
of $\proof$ ensures that triple $(\vec{\atv{}}, \msg, \proof)$ discloses no
information about the user  beyond what's revealed by $\vec{\atv{}}$.

\paragraph{Verification.} Verify is invoked by a verifier to check the validity
of credential presentation $(\vec{\atv{}}, \msg, \proof)$.  If the presentation
has already been produced by a valid present call, then $\F{\ATS}$ accepts and
returns true.  If it has already been flagged as invalid, then $\F{\ATS}$
rejects and returns false.  Finally, if the presentation has never been
received before, then $\F{\ATS}$ verifies whether the presentation could have
been created by corrupt users and issuers.  That is, $\F{\ATS}$ checks if there
are corrupt users that hold all the disclosed attribute values, and if not,
then $\F{\ATS}$ also checks if there are corrupt issuers that can produce valid
credentials for the attribute values not held by corrupt users.  In these
checks succeed, then $\F{\ATS}$ refers to $\simu$ and returns what $\simu$
returns; otherwise, it records the presentation, flags it as invalid, and
outputs false.

\paragraph{Register}. On input register from the user, send
$(register, user)$ to the registration authority and wait for its
confirmation.  Upon receiving the confirmation, do the following:

\begin{enumerate}

\item If $\creds[\user] \not\rightarrow \nil$, then abort; else continue.  (a
user can be registered only once})

\item Send $(register, \user)$ to \simu and wait for the go-ahead; on
  receiving the go-ahead continue.

\item Set $\creds[\user] \leftarrow [\varnothing, \ldots, \varnothing]$
  ($\creds[\user]$ is an array of length $N$, and $\creds[\user][i]$ is the
  $i^{th}$ element).

\item Return $\sf{register}end$ to the user .

\end{enumerate}

\paragraph{Issuance}. On input $(\issuecred, \att{i},\atv{i})$ from the user,
with $\atv{i}$ being value of attribute $\att{i}$, send $(\issuecred, \user,
\atv{i})$ to issuer $\issuer{i}$ associated with $\att{i}$ and wait for its
confirmation.  Upon confirmation, do:

\begin{enumerate}

  \item If issuer $\issuer{i}$ is honest, then:

    \begin{enumerate}

    \item If $\creds[\user] \rightarrow \nil$, then abort; else continue. (An
    honest issuer will only issue credentials to a registered user.)}

\item If $\creds[\user][i] \neq \varnothing$, then abort; else continue.  (An
honest issuer will issue at most one valid credential per user.)}

\item Send $(\issue, \user, \att{i}, \atv{i})$ to \simu and wait for the
  go-ahead; on receiving the go ahead, continue.

\end{enumerate}

\item Else:

  \begin{enumerate}

    \item Send $(\issuecred, \user, \att{i}, \atv{i})$ to $\simu$, and waits
      for its response.

    \item If $\simu$ sends $\accept$, then continue; else abort.

  \end{enumerate}

\item Set $\creds[\user][i] \leftarrow \creds[\user][i] \cup \{\atv{i}\}$.

\item Return $(\issueend, \att{i}, \atv{i})$ to \user and $(\issueend, \user,
\atv{i})$ to $\issuer{i}$.  \end{enumerate}

\paragraph{Aggregate}. On input $(aggregate, \vec{\atv{}})$ from the
user , where $\vec{\atv{}} = (\atv{1}, ..., \atv{N})$ and $\atv{i} \in
\{\varnothing\} \cup \{0, 1\}^*$, do:

\begin{enumerate}

\item If $\creds[\user] = \nil$, then abort; else continue.  (Unregistered
users cannot successfully aggregate credentials.)}

\item If $\forall i \in [1 \isep N]: \atv{i} = \varnothing \vee \atv{i} \in
  \creds[\user][i]$, then: \begin{enumerate} \item Randomly generate
    $\signature$.

  \item Set $\aggcreds[\signature] \leftarrow \langle \user,
  \vec{\atv{}}\rangle$.  \item Return $(\aggend, \signature)$ to the user .
(Users who hold non-empty attribute values successfully aggregate the
corresponding credentials.)}

\end{enumerate}

\item Else abort.

\end{enumerate}

\paragraph{Present.} On input $(\present, \vec{\atv{}}, \signature, \rand)$ from
the user , where $\vec{\atv{}} = (\atv{1}, ..., \atv{N})$ s.t. $\atv{i} =
\varnothing$ if $\atv{i}$ is not disclosed, do:

\begin{enumerate}

  \item If $\aggcreds[\signature] \rightarrow \langle \user,
    \vec{\atv{}}\rangle$, then:

    \begin{enumerate}

      \item Randomly generate $\proof$.  \item Set $\proven[\vec{\atv{}}, \msg,
        \proof] \leftarrow \valid$.

    \item Return $(\presentend, \proof)$ to the user  and $(\presentend,
    \vec{\atv{}}, \msg, \proof)$ to $\simu$.  (Users who successfully
  aggregated their credentials successfully sign messages.)}

\end{enumerate}

\item Else if $\creds[\user] = \nil$, then abort; else continue. (Unregistered
users cannot successfully sign messages.)}

\item If $\forall i \in [1\isep N]: \atv{i} = \varnothing \vee \atv{i} \in
  \creds[\user][i]$, then: \begin{enumerate}

    \item Randomly generate $\proof$.

  \item Set $\proven[\vec{\atv{}}, \msg, \proof] \leftarrow \valid$.  \item
  Return $(\presentend,  \proof)$ to the user  and $(\presentend, \vec{\atv{}},
\msg, \proof)$ to $\simu$. (Users with disclosed attribute values successfully
sign messages even if they did not call \sf{aggregate} before.)}
\end{enumerate}

\item Else abort.

\end{enumerate}

\paragraph{Verification.} Upon receiving message $(\verif, \vec{\atv{}}, \msg, \proof)$
from a verifier, do:

\begin{enumerate}

  \item If $\proven[\vec{\atv{}}, \msg, \proof] \rightarrow \invalid$, then
    abort.

  \item Else if $\proven[\vec{\atv{}}, \msg, \proof] \rightarrow \valid$, then:

    \begin{enumerate}

    \item Return $(\verifend, 1)$ to the verifier.  (Completeness.)}
            
  \item Send $(\verifend, \verifr, \vec{\atv{}}, \msg, \proof, \allowbreak
1)$ to \simu.  \end{enumerate}

\item Else if $\proven[\vec{\atv{}}, \msg, \proof] \rightarrow \nil$, then:

  \begin{enumerate}

    \item If there exists a dishonest user $\user'$ with attribute values
      $\vec{\atvprime{}} = (\atvprime{1}, \ldots, \atvprime{N})$ s.t.  $\forall
      \atv{i} \neq \varnothing$, we either have $\atv{i} = \atvprime{i}$, or
      have $\atv{i} \neq \atvprime{i}$ but issuer $\issuer{i}$ is dishonest,
      then query \simu with $(\verif, \vec{\atv{}}, \msg,\proof)$. On receiving
      the response $(\verifend, b)$, do: \begin{enumerate} \item Set
        $\proven[\vec{\atv{}}, \msg, \proof] \leftarrow b$.

        \item Return $(\verifend, b)$ to $\verifr$.

      \end{enumerate}

    \item Else:

      \begin{enumerate}

        \item Set $\proven[\vec{\atv{}}, \msg, \proof] \leftarrow \invalid$.

      \item Return $(\verifend, 0)$ to $\verifr$. (Unforgeability.)}

      \item Send $(\verifend, \verifr, \vec{\atv{}}, \msg, \proof,\allowbreak
        0)$ to \simu.

    \end{enumerate}

\end{enumerate}

\end{enumerate}

We assume there are $N$ issuers in the system, so the system supports $N$
attributes.  When $\F{\ATS}$ is called through aggregate, present, or
verify, it receives vector $\vec{\atv{}} = (\atv{1}, \ldots, \atv{N})$  such
that $\atv{i} = \varnothing$ or $\atv{i} \in \{0, 1\}^*$.  We remark that
$\atv{i} = \varnothing$ indicates that the value of attribute $\attr_{i}$ will
be hidden during credential presentation, whereas $\atv{i} \neq  \varnothing$
indicates that value is disclosed.

Let  $\textsc{Exec}_{\Pi,\adv}(\secparam)$ denote the outputs of the honest
parties and~$\adv$, during an execution of~$\Pi$ in the real world and
$\textsc{Exec}_{\F{\ATS},\simu}(\secparam)$ denote the outputs of the honest
parties and~$\simu$, interacting with~$\F{\ATS}$ in the ideal world.

\begin{definition} A protocol~$\Pi$ securely realises ideal
  functionality~$\F{\ATS}$ in the standalone model if, for every probabilistic
  polynomial time (PPT) adversary~$\adv$, there exists a PPT simulator~$\simu$
  satisfying $\textsc{Exec}_{\Pi,\adv}(\secparam) \approx_c
\textsc{Exec}_{\F{\ATS},\simu}(\secparam)$, where $\secparam$ is the security
parameter, and $\approx_c$ denotes computational indistinguishability.
\end{definition}

\subsection{Proof of Theorem \ref{thm:agg-creds-bb-security}}
\label{sec:proof-1}

\begin{proof}

  In the following, we show that the view of a real-world adversary $\adv$
  during an execution of $\Pi$ is indistinguishable from the view of an
  ideal-world simulator $\simu$ that interacts with system participants through
  ideal functionality $\F{\ATS}$.  We define a sequence of indistinguishable
  games that start with Game 0 representing the real world, and end with Game 8
  representing the ideal world.

  \paragraph{Game 0.} This is the real-world execution of protocol $\Pi$.

  \paragraph{Game 1.} The simulator $\simu$ executes real-world protocol
  $\Pi$ on behalf of the honest parties.  That is, $\simu$ is provided with the
  inputs (including secrets) of the honest parties, and
  interacts with corrupt parties, through the adversary, in the real-world,
  following $\Pi$'s specifications, and returns the corresponding output.  Game
  1 is by construction indistinguishable from Game 0.

  \paragraph{Game 2.} We introduce dummy functionality $\F{\ATS}'$ that
  forwards the inputs from the honest parties to $\simu$ and returns the
  respective outputs.  This modification is syntactic.

  \paragraph{Game 3.} In this game, $\F{\ATS}'$ handles register calls from
  honest users following the description of ideal functionality $\F{\ATS}$.
  $\simu$ first assigns each honest user a key pair $(\sk, \pk)$.  Now, when an
  honest user $\user$ calls $\F{\ATS}'$ to register, $\simu$ learns the user's
  identity.  Given this information, $\simu$ retrieves the corresponding key
  pair $(\sk,\pk)$ and contacts registration authority $\RA$ to register
  $\user$.  $\simu$ receives a unique tag $\tau$ and a signature $\psi
  \leftarrow \DSsign(\rask, \pk \| \tau)$ and stores entry $\langle \user,
  \tau, \pk, \psi\rangle $.

  Since $\RA$ is honest, they will register users only once and reject
  duplicate requests, same as $\F{\ATS}'$.  In other words, when $\F{\ATS}'$
  accepts a registration request, so will $\RA$ in the real world, and vice
  versa.  Hence, this game is indistinguishable from Game 2.

  \paragraph{Game 4.} In this game, $\F{\ATS}'$ also handles credential
  issuance calls from honest users following the description of ideal
  functionality $\F{\ATS}$.  When  $\F{\ATS}'$ receives issuance request
  $(\issue, \att{i}, \atv{i})$ from honest user $\user$, $\simu$ learns the
  triple $(\user, \att{i}, \atv{i})$.  To simulate the corresponding credential
  issuance for $\user$ in the real world, $\simu$ retrieves $(\sk, \pk)$ and
  $(\tau, \psi)$.  If the pair $(\tau, \psi)$ for $\user$ does not exist (i.e.,
  $\user$ did not call register before), then \simu aborts; otherwise, it
  honestly executes $\Pi$ with issuer $\issuer{i}$.  There are two cases to
  consider.

  \begin{enumerate}

    \item $\issuer{i}$ is honest.  If this is the first time $\simu$ simulates
      $\user$ for $\issuer{i}$, then as per $\Pi$'s specification,  $\simu$
      will receive signature $\signature_i \leftarrow \ATSsign(\isk{i}, \tau,
      \atv{i})$, and record entry $\langle \user, \tau, \atv{i}, \signature_i
      \rangle$.  Otherwise, $\issuer{i}$ will reject the credential issuance
      request.

    \item $\issuer{i}$ is corrupt.  If the execution of $\Pi$ terminates with
      $\simu$ receiving signature $\signature_i : \ATSverify(\ipk{i}, \tau,
      \atv{i}, \signature_i) = 1$, then $\simu$ records entry
      $\langle \user, \tau, \atv{i}, \signature_i \rangle$, and sends \accept.

  \end{enumerate}

  We recall that, in the real world, an honest issuer rejects a credential
  issuance request from a registered user if a preceding request from the same
  user has already been processed.  Otherwise, the honest issuer accepts and
  issues the credential.  It is easy to see that, in this case, the output of
  the issuer matches that of $\F{\ATS}'$ if the latter processes the credential
  request instead.  We also recall that when $\F{\ATS}'$ receives a  credential
  issuance request involving a dishonest issuer, $\F{\ATS}'$ follows $\simu$
  instructions.  That is, if $\simu$ sends accept, then $\F{\ATS}'$ accepts
  the request and updates the credential map accordingly; otherwise,
  $\F{\ATS}'$ just rejects.  Since $\simu$ sends accept only if the dishonest
  issuer outputs a valid aggregate signature, then the output of the dishonest
  issuer in the real world is equivalent to the output of $\F{\ATS}'$.
  Therefore, this game is indistinguishable from Game 4 as long as $\simu$ does
  not abort.  Since $\user$ is honest, she will always register first before
  requesting credentials from the issuers, and hence, the abort event will
  never occur.

  \paragraph{Game 5.} During this game, $\F{\ATS}'$ executes aggregate requests
  from honest users following the description of $\F{\ATS}$.  This game is
  indistinguishable from Game 4, as credential aggregation, in the real world,
  is a local operation during which adversary $\adv$ observes nothing, and
  therefore, $\simu$ need not do anything.

  \paragraph{Game 6.} In this game, $\F{\ATS}'$ takes charge of signature
  requests from honest users, following the specifications of $\F{\ATS}$, and
  $\simu$ successfully simulates their executions of $\Pi$ in the real world
  thanks to the simulatability of the signature of knowledge $\sok$.  Namely,
  given message $(\presentend, \vec{\atv{}}, \msg, \proof)$ from $\F{\ATS}'$
  for $\vec{\atv{}} = (\atv{1}, \ldots, \atv{N})$, $\simu$ defines  $\inst =
  (\pk, (\ipk{i}, \atv{i})_{i \in \ind})$ such that $\forall i \in \ind:
  \atv{i} \neq \varnothing$, and outputs $\proof^{\star} \leftarrow
  \soksimsign(\td, \inst, \msg, \crs)$, which gives $\sokverify(\inst,
  \proof^{\star}, \crs) = 1$.

  \paragraph{Game 7.} In this game, $\F{\ATS}'$ processes verification requests
  from honest verifiers following the description of $\F{\ATS}$, whereas
  $\simu$ simulates the honest verifiers in the real world.  Namely, upon
  receiving a signature in the real world, $\simu$ translates the signature
  into calls to $\F{\ATS}'$, and instead of returning the output of
  $\sokverify$, it returns the verification result received from $\F{\ATS}'$.
  In the following, we describe $\simu$'s simulation.

  \paragraph{Registration.} For each registration request $\pk$ that a corrupt
  user sends, in the real world, to the honest registration authority, $\simu$ 
  (acting as $\RA$) first checks if the user knows the secret key
  $\sk$ underlying $\pk$.  If not, then $\simu$ stops.  Otherwise, it executes
  the rest of the registration protocol, records entry $\langle\user, \tau,
  \pk, \psi \rangle$, and then (acting as a user) transmits the corresponding
  registration request $register$ to $\F{\ATS}'$ which processes it in a
  similar fashion to $\F{\ATS}$.

  As per the specifications of $\Pi$, the simulated $\RA$ in the real world
  will not accept a registration request that will be rejected by $\F{\ATS}'$,
  and vice versa.  Therefore the output of $\Pi$ and the output of $\F{\ATS}'$
  when processing registration requests is indistinguishable.

  \paragraph{Credential Issuance.} For each issuance request $(\tau, \pk,
  \psi, \atv{i})$ that a corrupt user $\user$ sends, in the real world, to an
  honest issuer $\issuer{i}$ (simulated by $\simu$), $\simu$ (acting as
  $\user$) sends $(\issue, \att{i}, \atv{i})$ to $\F{\ATS}'$.  We recall that
  if $\user$ is not registered or if this is not the first time that $\user$
  requests a credential, $\F{\ATS}'$ rejects.  Otherwise, $\F{\ATS}'$ issues
  the credential.  Similarly, in the real world, the simulated $\issuer{i}$
  returns $\signature_i \leftarrow \ATSsign(\isk{i}, \tau, \atv{i})$ if the
  same conditions are met.  Hence, the output of $\Pi$ and the output of
  $\F{\ATS}'$ when processing credential issuance is indistinguishable.  After
  each successful credential issuance, $\simu$ stores entry $\langle \user,
  \tau, \att{i}, \atv{i}, \signature_i\rangle$.

  \paragraph{Signature Verification.} For each signature $(\vec{\atv{}} =
  (\atv{1}, ..., \atv{N}), \msg, \proof)$ that a user $\user$ presents, in the
  real world, to an honest verifier $V$, $\simu$ (simulating $V$) calls $b
  \leftarrow \sokverify(\crs, \vec{\atv{}}, \msg, \proof)$.  We recall that
  $\atv{i} = \varnothing$ if the attribute value is not disclosed.

  If $b  = 0$,  $\simu$ sends $(\msg, \vec{\atv{}} = (\atv{1}, ...,
  \atv{N}), \proof)$ to $\F{\ATS}'$, which returns $(\verifend, b^\star)$.
  Two scenarios present themselves. Either rhere exists a dishonest user
  with attribute values $\vec{\atvprime{}} = (\atvprime{1}, ..., \atvprime{N})$,
   $\forall \atv{i} \neq \varnothing$, we either have $\atv{i} =
  \atvprime{i}$, or have $\atv{i} \neq \atvprime{i}$ but issuer $\issuer{i}$ is
  dishonest.  Or $\forall i: \atv{i} \neq \varnothing$, issuers
  $\issuer{i}$ and users holding attribute values $\atv{i}$ are honest.

  In case of scenario one, the output of $\F{\ATS}'$ is defined by
  $\simu$, which instructs $\F{\ATS}'$ to return $(\verifend, 0)$.  In
  case of scenario two, $\F{\ATS}'$ outputs $(\verifend, 1)$.
  However, thanks to the correctness of protocol $\Pi$ this scenario never
  occurs.  That is, given a signature of knowledge $\proof$ that was computed
  honestly,  $\sokverify$ will never output $0$.  Thus, the output of
  $\F{\ATS}'$ matches that of $\Pi$.

  If $b = 1$, then thanks to the extractability property of the
  signature of knowledge $\sok$, $\simu$ extracts witness $\witness = (\tau,
  \sk, \pk, \psi, \signature)$  and then executes the following steps:

  \begin{itemize}

    \item If there is no recorded entry $\langle \star, \tau, \pk,
      \star\rangle$, then $\simu$ aborts.

    \item Else if there exists $1 \leq i \leq N$ such that $\atv{i} \neq
      \varnothing$ and $\issuer{i}$ is honest, and there is no recorded entry
      $\langle \star, \tau, \atv{i}, \star \rangle$, then $\simu$ aborts.

    \item Otherwise, $\simu$ sends to $\F{\ATS}'$ signature request $(\sign,
      \vec{\atv{}}, \signature, \msg)$.

      \begin{itemize}

        \item If $\F{\ATS}'$ aborts, then $\simu$ sends request $(\verif,
          V,\vec{\atv{}}, \proof)$ to $\F{\ATS}'$.

        \item Otherwise, $\F{\ATS}'$ returns $(\presentend,  \proof^{\star})$,
          and  $\simu$ sends request $(\verif, V, \vec{\atv{}},
          \proof^{\star})$ to $\F{\ATS}'$.

      \end{itemize}

      As a result, $\simu$ receives $(\verifend, b)$ and outputs $b$ in
      turn.

  \end{itemize}

  In what follows, we show that if $\simu$ does not abort, then $b =
  1$.

  We call the two abort events $\abort{1}$ and $\abort{2}$.  If $\abort{1}$
  does not occur, then this implies that a registration request for $\user$ was
  successfully processed by $\simu$ and therefore $\user$ was successfully
  registered by $\F{\ATS}$.  If $\abort{2}$ does not occur then this implies
  one of two things: 
  
case one $\forall i : \atv{i} \neq \varnothing$, there
  exists an entry $\langle \star, \tau, \atv{i}, \star \rangle$; 
  
case two 
  $\exists j: \atv{j} \neq \varnothing$ such that there is no recorded entry
  $\langle \star, \tau, \atv{j}, \star \rangle$ but issuer $\issuer{j}$ is
  dishonest.  
  
case one implies that $\simu$ issued credentials for $\user$ for
  the relevant attributes, and so did $\F{\ATS}'$.  Therefore, during the
  $\present$ call, $\F{\ATS}'$ will not abort and will return $(\presentend,
  \proof^{\star})$, and the subsequent call $(\verif, V, \vec{\atv{}},
  \proof^{\star})$ will result in $\F{\ATS}'$ returning $(\verifend, 1)$.
  This matches the output of the verification in the real world.  
  
case two 
  implies that $\user$ is in possession of a credential that was issued by a
  dishonest issuer and never recorded by either $\simu$ or $\F{\ATS}'$.  In
  this case,  $\F'{}$ aborts during $\present$, and when $\simu$ calls
  $\F{\ATS}'$ with  $(\verif, V, \vec{\atv{}}, \proof)$, the output of
  $\F{\ATS}'$ is dictated by $\simu$.  $\simu$ sets this output to $1$ to
  match the output of the verification in the real world.

  Therefore, this game is indistinguishable from Game 6 as long as $\abort{1}$
  and $\abort{2}$ do not occur.

  Now we show that the probability that one of these two events occur is
  negligible.

  $\abort{1}$ occurs if there is no recorded entry $\langle \star, \tau, \pk,
  \star \rangle$ in the system.  Again due to extraction of the verifying
  signature of knowledge $\proof$, we can extract witness $\witness = (\tau,
  \sk, \pk, \psi, \signature)$.  This means that $\user$ has knowledge of a
  valid $\psi$ such that $1 \leftarrow \DSverify(\pk, \pk \| \tau, \psi)$ with
  no corresponding entry, contradicting the existential unforgeability of the
  digital signature scheme $\DS$.

  $\abort{2}$ occurs if $\exists j \in [1\isep N]$ such that $\atv{j} \neq
  \varnothing$ and $\issuer{j}$ is honest, and there is no recorded entry
  $\langle \star, \tau, \atv{j}, \star \rangle$.  Again, due to extraction of
  the correctly verifying signature of knowledge $\proof$, we can extract a
  valid witness $\witness = (\tau, \sk, \pk, \psi, \signature)$.  This means
  that $\signature$ satisfies $\ATSverifyagg((\ipk{i})_{i\in \ind}, \tau,
  (\atv{i})_{i\in\ind}, \signature)$, where $\ind$ identifies the indices s.t.
  $\atv{i} \neq \varnothing$ and $j \in \ind$.  Given that there is no recorded
  entry $\langle \star, \tau, \atv{j}, \star \rangle$, this implies that the
  honest issuer $\issuer{j}$ has never issued a signature $\signature_j$ for
  pair $(\tau, \atv{j})$, which contradicts the existentially unforgeability
  under tag-based aggregation of $\ATS$.

\paragraph{Game 8.} This is the ideal world.  This game is indistinguishable
from Game 7 as functionality $\F{\ATS}'$ is identical to ideal functionality
$\F{\ATS}$.  \end{proof}

\subsection{Proof of Theorem \ref{thm:tag-uf}}\label{app:tag-uf}

\begin{proof} Let $\pk = \{\pk_1, ..., \pk_N\}$ be the set of challenge
  public keys from the existential unforgeability under tag-based aggregation
  experiment.  Assume there exists an adversary $\adv$ that succeeds in the
  unforgeability experiment by producing an aggregate signature
  $(\signatureone, \signaturetwo)$ on a set of messages $(\mssg_i^*)_{i \in
  \ind \subseteq [1 \isep N]}$ under a tag $\tau^*$ such that the following
  holds: \begin{equation*} e(\signaturetwo,\ggen) = e(\hashtwo(\tau^*),
    \signatureone)\cdot \prod_{i \in \ind} e(\hashone(\mssg_i^*), \pk_i).
  \end{equation*} In addition, there exists an index $j \in \ind$ such that the
  adversary does not control the secret key corresponding to $\pk_j$, and a
  tag-based aggregate signature has not been issued by the signature oracle on
  triple $(\pk_j, \tau^*, \mssg_j^*)$.  We distinguish between two types of
  forgeries.

  \begin{description}

    \item[Type 1] \adv has never issued a signature request for $(\pk_j,
      \star, \mssg^*_j)$.

    \item[Type 2]  \adv has issued signature requests for $(\pk_j, \tau,
      \mssg^*_j)$ but $\forall \tau: \tau \neq \tau^*$.

  \end{description}

  In the following, we show that: (1) if adversary $\adv$ outputs a Type 1
  forgery, then there exists an adversary $\advB$ that leverages  $\adv$ to
  break BCDH, and (2) if $\adv$ outputs a Type 2 forgery, then there exists an
  adversary $\advB$ that leverages $\adv$ to break m-BCDH.

  \textbf{Type 1 Forgery.} Let $\advB$ be an adversary whose goal is to break
  BCDH, and hence, is given tuple $(\gen, X, Y, Z, \ggen, \widehat{X},
  \widehat{Y}, \widehat{Z})= (\gen, \gen^x, \gen^y, \gen^z, \ggen, \ggen^x,
  \ggen^y, \ggen^z)$.  To compute $\gen^{xy}$ and break BCDH, $\advB$ leverages
  $\adv$ and the random oracle as depicted below.

  \paragraph{Random Oracle Programming.} $\advB$ initialises two lists
  $\ListOne$ and $\ListTwo$, and then computes $\hashone$ and $\hashtwo$ as
  follows:

  \begin{itemize}

    \item $\hashone$: For each query $\mssg$, check if there exists an entry
      $\langle \mssg, \star, \hashone(\mssg) \rangle \in \ListOne$, and if so,
      return $\hashone(\mssg)$.  If not, then do the following:

      \begin{itemize}

        \item With probability $p_\hashone$, generate a random element $\rho
          \leftarrow \Zp$, set $\hashone(\mssg) = \gen^{\rho}$, and store entry
          $\langle \mssg, \rho, \hashone(\mssg)\rangle$ in $\ListOne$;

        \item With probability $1 - p_\hashone$, set $\hashone(\mssg) = Y$ and
          store entry $\langle \mssg, \perp, \hashone(\mssg)\rangle$ in
          $\ListOne$.

      \end{itemize}

    \item $\hashtwo$: For each query $\tau$, check if there exists an entry
      $\langle \tau, \star, \hashtwo(\tau) \rangle \in \ListTwo$, and if so,
      return $\hashtwo(\tau)$.  If not, then select a random element $\eta
      \leftarrow \Zp$, set $\hashtwo(\tau) = \gen^{\eta}$, and store entry
      $\langle \tau, \eta, \hashtwo(\tau)\rangle$ in $\ListTwo$.

  \end{itemize}

  \paragraph{\Corrupt Simulation.} Without loss of generality, we assume that
  $\adv$ corrupts all public keys except for one.  $\advB$ randomly selects $j
  \in [1\isep N]$ and sets $\pk_j = \widehat{X}$, and for all $k \neq j$
  randomly selects $x_k \in \Zp^*$ and has $\sk_k = x_k$ and $\pk_k =
  \widehat{X}_k = \ggen^{x_k}$.  If $\adv$ sends a corruption request for
  $\pk_i \neq \pk_j$, then $\advB$ returns $\sk_i = x_i$; otherwise, $\advB$
  aborts the experiment.

  \paragraph{\Sign Simulation.} On receiving a signature request $(\pk_i,
  \tau, \mssg)$, $\advB$ proceeds as follows.  If $i \neq j$, then $\advB$
  outputs $\signature \leftarrow \ATSsign(\tau, \sk, \mssg)$.  If $i = j$,
  then $\advB$ checks if $\hashone(\mssg) = Y$. If so, then $\advB$ aborts,
  otherwise,  $\advB$ retrieves entry $\langle \mssg, \rho, \hashone(\mssg) =
  \gen^{\rho}\rangle$ and computes $\signature = (\signatureone, \signaturetwo)
  = (\ggen^a, X^{\rho}\hashtwo(\tau)^a)$.  Note that $\signature$ is a valid
  signature on pair $(\tau, \mssg)$ relative to public key $\pk_j =
  \widehat{X} = \ggen^x$; in fact,  $\signature = (\ggen^a,
  \hashone(\mssg)^{x}\hashtwo(\tau)^a)$.

  \paragraph{Forgery.} After receiving $\adv$'s forgery $(\ind,
  (\pk_i)_{i\in\ind}, \tau^*,  (\mssg^*_i)_{i \in \ind}, \signature^*)$,
  $\advB$ proceeds as follows.  $\advB$ checks if $\hashone(\mssg^*_j) \neq Y$,
  and if so, aborts.  If $\hashone(\mssg^*_j) = Y$, then $\advB$ parses
  $\signature^*$ as $(\signatureone^*, \signaturetwo^*)$, retrieves entry
  $\langle\tau^*, \eta^*, \hashtwo(\tau^*) \rangle$ from $\ListTwo$, and
  computes $\Delta = \frac{\signaturetwo^*}{\prod_{i \in \ind, i \neq
  j}\hashone(\mssg^*_i)^{x_i}}$ followed by $\Gamma = \frac{e(\Delta,
\widehat{Z})}{e(Z^{\eta^*},\signatureone^*)}$.

We now prove that $\Gamma = e(\gen, \ggen)^{xyz}$.  Note that  $\signature^* =
(\signatureone^*, \signaturetwo^*)$ verifies the following equality:
$e(\signaturetwo^*, \ggen) = e(\hashtwo(\tau^*), \signatureone^*) \cdot
\prod_{i\in \ind}e(\hashone(\mssg^*_i), \widehat{X}_i).$

Therefore, $e(\Delta, \ggen) = e(\hashtwo(\tau^*), \signatureone^*) \cdot
e(\hashone(\mssg^*_j), \widehat{X}_j) = e(\gen^{\eta^*}, \signatureone^*)e(Y,
\widehat{X})$.  It follows that we have $e(Y, \widehat{X}) = \frac{e(\Delta,
\ggen)}{e(\gen^{\eta^*}, \signatureone^*)}$ and $e(Y, \widehat{X})^z =
\frac{e(\Delta, \widehat{Z})}{e(Z^{\eta^*}, \signatureone^*)} = \Gamma$.

To summarize, $\advB$ aborts in the following cases:

\begin{itemize}

  \item  If $\adv$ sends a corruption request for $\pk_i = \pk_j$. As we
    assume that $\adv$ corrupts all public keys except one, and we have $j$
    chosen uniformly at random, the probability that $\advB$ \emph{succeeds}
    here is $1/N$.

  \item If, for the forged message under the uncorrupted key $\pk_j$, we have
    the case that $\hashone(\mssg^*_j) \neq Y$, then $\advB$ cannot extract
    $e(\gen,\ggen)^{xyz}$ as required. This occurs with probability
    $p_{\hashone}$.  Hence, the probability that $\advB$ \emph{succeeds} here
    is $1 - p_{\hashone}$.

  \item If any signing query $(\pk_j,\tau,\mssg)$ to the signing oracle is for
    a message such that $\hashone(\mssg) = Y$, then $\advB$ cannot produce a
    signature and must abort. The probability that all signing queries avoid
    this is $p_{\hashone}^{\sigqueries}$, where $\sigqueries$ is the number of
    signing queries made with $\pk_j$.

\end{itemize}

The probability of success is then $\frac{1}{N}\cdot (1-p_{\hashone}) \cdot
p_{\hashone}^{\sigqueries}$, and so the probability of aborting in the case of
Type 1 forgeries is given as follows: \[ \Pr[\text{Type 1 abort}] = 1 - \left(
\frac{1}{N}\cdot (1-p_{\hashone}) \cdot p_{\hashone}^{\sigqueries} \right). \]

\textbf{Type 2 Forgery.} Let $\advB$ be an adversary whose goal is to break
m-BCDH, and therefore, has as input $(\gen, Y, \ggen, \widehat{X},
\widehat{Y})= (\gen, \gen^{1/y}, \ggen, \ggen^x, \ggen^y)$.  To break m-BCDH,
$\advB$ simulates the unforgeability experiment to $\adv$ as follows.

\paragraph{Random Oracle Programming.} $\advB$ initializes two lists $\ListOne$
and $\ListTwo$, and then computes $\hashone$ and $\hashtwo$ as described below:

\begin{itemize}

  \item $\hashone$: For each query $\mssg$, check if there exists an entry
    $\langle \mssg, \star, \hashone(\mssg) \rangle \in \ListOne$, and if so,
    return $\hashone(\mssg)$.  If not, then generate a random element $\rho
    \leftarrow \Zp$, set $\hashone(\mssg) = \gen^{\rho}$, and store entry
    $\langle \mssg, \rho, \hashone(\mssg)\rangle$ in $\ListOne$.

  \item $\hashtwo$: For each query $\tau$, check if there exists an entry
    $\langle \tau, \star, \hashtwo(\tau) \rangle \in \ListTwo$, and if so,
    return $\hashtwo(\tau)$.  If not, then first generate a random element
    $\eta \in \Zp$, then with probability $p_\hashtwo$, set $\hashtwo(\tau) =
    Y^{\eta}$, and with probability $1 - p_\hashtwo$, set $\hashtwo(\tau) =
    \gen^{\eta}$,  finally, store entry $\langle \tau, \eta,
    \hashtwo(\tau)\rangle$ in $\ListTwo$.

\end{itemize}

\paragraph{\Corrupt Simulation.} Without loss of generality, we assume that
$\adv$ corrupts all public keys except for one.  $\advB$ randomly selects $j
\in [1\isep N]$ and sets $\pk_j = \widehat{X}$, and for all $k \neq j$
randomly selects $x_k \in \Zp^*$ and has $\sk_k = x_k$ and $\pk_k =
\widehat{X}_k = \ggen^{x_k}$.  If $\adv$ sends a corruption request for $\pk_i
\neq \pk_j$, then $\advB$ returns $\sk_i = x_i$; otherwise, $\advB$ aborts
the experiment.

\paragraph{\Sign Simulation.} On receiving a signature request $(\pk_i, \tau,
\mssg)$, $\advB$ proceeds as follows.  If $i \neq j$, then $\advB$ outputs
$\signature \leftarrow \ATSsign(\tau, \sk, \mssg)$.  If $i = j$, then $\advB$
first retrieves entries $\langle \mssg, \rho, \hashone(\mssg)\rangle$ and
$\langle \tau, \eta, \hashtwo(\tau)\rangle$ from $\ListOne$ and $\ListTwo$
respectively.  Now if $\hashtwo(\tau) \neq \gen^{\eta}$ or $\eta = 0$, then
$\advB$ aborts.  Otherwise, $\advB$ randomly selects $a \in \Zp$, lets
$\signaturetwo = \gen^{a}$ and  computes $\signatureone =
(\ggen^a/\widehat{X}^{\rho})^{1/\eta}$.  Note that $\signature$ is a valid
signature on pair $(\tau, \mssg)$ relative to public key $\pk_j = \widehat{X}
= \ggen^x$; in fact, by construction, $e(\hashtwo(\tau), \signatureone) \cdot
e(\hashone(\mssg), \widehat{X}) = e(\signaturetwo, \gen)$.

\paragraph{Forgery.} After receiving $\adv$'s forgery $(\ind,
(\pk_i)_{i\in\ind}, \tau^*,  (\mssg^*_i)_{i \in \ind}, \signature^*)$, $\advB$
proceeds as follows.  $\advB$ first retrieves entries $\langle \mssg^*, \rho^*,
\hashone(\mssg^*) = \gen^{\rho^*}\rangle$ and $\langle \tau^*, \eta^*,
\hashtwo(\tau^*)\rangle$ from $\ListOne$ and $\ListTwo$ respectively.

$\advB$ next checks if $\hashtwo(\tau^*) \neq Y^{\eta^*}$ or $\rho^* = 0$, and
if so, aborts.  Otherwise, $\advB$ parses $\signature^*$ as $(\signatureone^*,
\signaturetwo^*)$, and computes $\Delta = \frac{\signaturetwo^*}{\prod_{i \in
\ind, i \neq j}\hashone(\mssg^*_i)^{x_i}}$ followed by $\Gamma =
\left(\frac{e(\Delta,
\widehat{Y})}{e(\gen^{\eta},\signatureone^*)}\right)^{1/\rho^*}$.

We show now that $\Gamma = e(\gen, \ggen)^{xy}$.  Recall that  $\signature^* =
(\signatureone^*, \signaturetwo^*)$ verifies the following equality:
$e(\signaturetwo^*, \ggen) = e(\hashtwo(\tau^*), \signatureone^*) \cdot
\prod_{i\in \ind}e(\hashone(\mssg^*_i), \widehat{X}_i).$ Therefore: $e(\Delta,
\ggen) = e(\hashtwo(\tau^*), \signatureone^*) \cdot e(\hashone(\mssg^*_j),
\widehat{X}_j) = e(Y^{\eta^*}, \signatureone^*) \cdot e(\gen^{\rho^*},
\widehat{X})$.

Hence, $e(\gen^{\rho^*}, \widehat{X})^y  = \frac{e(\Delta,
\widehat{Y})}{e(\gen^{\eta^*}, \signatureone^*)}$.  This implies that $e(\gen,
\widehat{X})^y = e(\gen, \ggen)^{xy} = \left(\frac{e(\Delta,
\widehat{Y})}{e(\gen^{\eta^*}, \signatureone^*)}\right)^{1/\rho^*}$, which
equals $\Gamma$.

In the case of Type 2 forgeries, $\advB$ aborts in the following cases:

\begin{itemize}

  \item If $\adv$ sends a corruption request for $\pk_i = \pk_j$. As we again
    assume $\adv$ corrupts all public keys except one, and $j$ is chosen
    uniformly at random, the probability that $\advB$ \emph{succeeds} here is
    again $1/N$.

  \item If, for the forged tag $\tau^*$, we have that $\hashtwo(\tau^*) \neq
    Y^{\eta^*}$, then $\advB$ cannot extract $e(\gen,\ggen)^{xy}$ as required.
    This occurs with probability $1-p_{\hashtwo}$.  Hence, the probability that
    $\advB$ \emph{succeeds} here is $p_{\hashtwo}$.

  \item If any signing query $(\pk_j,\tau,\mssg)$ to the signing oracle is for
    a tag such that $\hashtwo(\tau) \neq Y^{\eta}$, then $\advB$ cannot produce
    a signature and must abort. The probability that all signing queries avoid
    this is $p_{\hashtwo}^{\sigqueries}$, where $\sigqueries$ is the number of
    signing queries made with $\pk_j$.

\end{itemize}

The probability of success is then $\frac{1}{N}\cdot
p_{\hashtwo}^{\sigqueries+1}$, and so the probability of aborting in the case
of Type 2 forgeries is given as follows: \begin{equation*} \Pr[\text{Type 2
  abort}] = 1 - \left( \frac{1}{N}\cdot p_{\hashtwo}^{\sigqueries+1} \right).
\end{equation*}

\end{proof}
